\documentclass[a4,12pt]{article}
\usepackage[utf8]{inputenc}
\usepackage[T2A]{fontenc}
\usepackage[english, russian]{babel}

\sloppy % Все тексты в документе будут верстаться с менее строгими правилами
\usepackage{csquotes}
\usepackage{epigraph}
\usepackage{tcolorbox}
% Определяем окружение для Формата ввода
\newtcolorbox{inputformat}[1][]{colback=blue!5!white, colframe=blue!65!black,
    fonttitle=\bfseries, title=Формат ввода, #1}

% Определяем окружение для Формата вывода
\newtcolorbox{outputformat}[1][]{colback=green!5!white, colframe=green!65!black,
    fonttitle=\bfseries, title=Формат вывода, #1}

% % Определяем окружение для комментария
% \newtcolorbox{exercisecomment}[1][]{colback=yellow!5!white, colframe=yellow!80!black, fonttitle=\bfseries, title=Комментарий, #1}

% % Определяем окружение для замечания
% \newtcolorbox{exercisenote}[1][]{colback=red!5!white, colframe=red!80!black, fonttitle=\bfseries, title=Замечание, #1}

\newtcolorbox{exercisecomment}[1][]{colback=yellow!5!white, colframe=yellow!80!black, fonttitle=\bfseries, title=Комментарий, #1}

\newtcolorbox{exercisenote}[1][]{colback=red!5!white, colframe=red!80!black, fonttitle=\bfseries, title=Замечание, #1}

\newtcolorbox{exercisehint}[1][]{colback=gray!5!white, colframe=gray!60!black, fonttitle=\bfseries, title=Подсказка, #1}

\usepackage{geometry}
%\geometry{papersize={20.3 cm,25.4 cm}}
\geometry{left=2cm}
\geometry{right=2cm}
\geometry{top=2cm}
\geometry{bottom=2cm}


\usepackage{amsmath, amsfonts, amsthm, amssymb, amsopn, amscd}
\usepackage{enumerate}
\usepackage{enumitem}
\usepackage[mathscr]{eucal}

\usepackage{hyperref}
\hypersetup{unicode=true,final=true,colorlinks=true}

\theoremstyle{remark}
%\newtheorem{exercises}{Упражнения}
\newtheorem{exercise}{\textbf{Упражнение}}[section]
\renewcommand{\theexercise}{\textbf{\arabic{exercise}}}
%\renewcommand{\theexercise}{\textbf{\#\arabic{exercise}}}


\title{Листок 08. Циклы в Python}

\author{А.Н. Лата}

\begin{document} 

%\maketitle

\section*{\centering Листок 08. Циклы в Python}

\begin{exercisenote}[title=Замечания]
\begin{itemize}
    \item для решение приведенных ниже упражнений не требуется создавать (определять) пользовательские функции и использовать сторонние библиотеки.
    \item при решении упражнений полученные результаты выведите на экран с помощью функции {\color{blue}{print()}}.
    \item позаботесь о том, чтобы выводимый на экран результат был снабжен информацией о нем, там где это необходимо.
    \item не забывайте писать комментарии к вашему коду
\end{itemize}
\end{exercisenote}

%Выведите означает выведите на экран

\centering{\textbf{\Large Упражнения}}
\begin{enumerate}
    \subsection*{Начальный уровень}
    \item Выведите числа от 1 до 10.
    \item Выведите все четные числа от 2 до 20.
    \item Выведите все числа от 1 до 10 в обратном порядке.
    \item Выведите сумму чисел от 1 до 100.
    \item Выведите сумму всех нечетных чисел от 1 до 50.
    \item Выведите произведение чисел от 1 до 10.
    \item Выведите квадраты чисел от 1 до 10.
    \item Выведите таблицу умножения на 5.
    \item Выведите все буквы строки \texttt{'Python'} по одной в строке.

    \item Подсчитайте количество цифр в строке \texttt{'Python 3.14.2'}.
    \item Выведите символы строки \texttt{'Python'} в обратном порядке.

    % Исправление: добавлено уточнение «без использования списковых включений»,
    % чтобы задача не дублировала задачу №1 из раздела «Списковые включения».
    \item Создайте список, содержащий квадраты чисел от 1 до 10
    (без использования списковых включений).

    \subsection*{Средний уровень}
    \item Дана строка \texttt{'Ананас — вкусный фрукт, а ананас — тропический'}.
    Подсчитайте, сколько раз в строке встречается буква \texttt{'а'}
    в любом регистре (то есть учитывайте и \texttt{'а'}, и \texttt{'А'}).
    Выведите раздельно: количество строчных \texttt{'а'}, количество
    заглавных \texttt{'А'} и их сумму.
    \item Выведите все числа от 1 до 100, которые делятся на 7.
    \item Выведите все числа от 1 до 100, пропуская числа, которые делятся на 5.
    \item Найдите все числа от 1 до 50, которые делятся на 3 и на 5 одновременно.
    \item Создайте список из всех делителей числа 36.
    \item Создайте список с квадратами нечетных чисел от 1 до 20.
    \item Найдите сумму квадратов чисел от 1 до 20.
    \item Найдите сумму первых 10 натуральных чисел, кратных 3.

    % Исправление: добавлена подсказка об операциях // и %, так как
    % извлечение цифр числа — нетривиальный для начинающих прием.
    \item Найдите все двузначные числа, сумма цифр которых равна 8.

    \begin{exercisehint}
        Для числа \texttt{n} его десятки можно получить с помощью \texttt{n~//~10},
        а единицы — с помощью \texttt{n~\%~10}.
    \end{exercisehint}

    % Исправление: добавлена конкретная строка и уточнен язык (английский),
    % так как наборы гласных в русском и английском различаются.
    \item Подсчитайте количество гласных и согласных букв в строке \texttt{'Hello World'}
    (гласные английского языка: \texttt{a, e, i, o, u}; пробелы не учитываются).

    % Исправление: добавлены конкретные числа для проверки результата;
    % добавлена подсказка об алгоритме, поскольку задача нетривиальна для начинающих.
    \item Найдите наибольший общий делитель чисел 48 и 36.

    \begin{exercisehint}
        Перебирайте числа от 1 до меньшего из двух заданных и проверяйте,
        делится ли на него каждое из чисел без остатка. Наибольший из таких
        делителей и будет ответом.
    \end{exercisehint}

    \item Выведите первые 10 чисел последовательности Фибоначчи.
    
    \begin{exercisehint}
        Последовательность Фибоначчи начинается с двух чисел 0 и 1,
    каждое следующее число равно сумме двух предыдущих:
    $0,\; 1,\; 1,\; 2,\; 3,\; 5,\; 8,\; 13,\; \ldots$

    Заведите две целочисленные переменные \texttt{a = 0} и \texttt{b = 1},
    хранящие два соседних числа последовательности.
    На каждом шаге цикла выводите \texttt{a}, затем обновляйте
    обе переменные одновременно: \texttt{a, b = b, a + b}.
    Одновременное присваивание гарантирует, что старое значение
    \texttt{a} не будет потеряно до вычисления суммы.
    \end{exercisehint}

    % Исправление: добавлено конкретное число для проверки результата.
    \item Проверьте, является ли число \texttt{12321} палиндромом.
    
    \begin{exercisehint}
    Число называется \emph{палиндромом}, если оно читается одинаково слева направо и справа налево. Например: $12321$, $9009$.
    \end{exercisehint}

    \item Подсчитайте сумму цифр числа 12345.
    \item Выведите каждый второй символ строки \texttt{'abcdef'}.
    \item Найдите сумму всех элементов списка \texttt{[1, 2, 3, 4, 5]}, не используя функцию \texttt{sum()}.
    \item Создайте новый список, содержащий только четные числа из списка \texttt{[1, 2, 3, 4, 5, 6, 7, 8, 9, 10]}.
    \item Вычислите произведение всех чисел в списке \texttt{[1, 3, 5, 7, 9]}.
    \item Вычислите сумму всех чисел в двумерном списке \texttt{[[1, 2], [3, 4], [5, 6]]}.
    \item Найдите сумму всех нечетных чисел в строке \texttt{'1a2b3c4d5e6f7g8h9i'}.
    \item Найдите факториал числа 10.
    \item Выведите все гласные буквы в строке \texttt{'Hello World'}.
    \item Подсчитайте, сколько раз встречается слово \texttt{'apple'} в списке \texttt{['apple',~'banana',~'apple',~'orange']}.
    \item Создайте список, содержащий строки длиной больше 3 символов из списка \texttt{['cat',~'elephant',~'dog',~'mouse']}.
    \item Выведите первый символ каждого слова в строке \texttt{'Hello World from Python'}.
    \item Составьте строку, состоящую из первых букв каждого слова в строке \texttt{'Python is great'}.

    % Исправление: заменена строка 'I love programming' на 'We love programming',
    % так как слово 'I' имеет длину 1 и срез word[:3] дает неоднородный результат.
    \item Создайте строку, состоящую из первых трех букв каждого слова в строке \texttt{'We love programming'}.

    % Исправление: заменена строка на готовый список, чтобы убрать неоднозначность
    % разбиения строки с запятыми и пробелами методом split().
    \item Сформируйте список, содержащий только те элементы, которые начинаются на букву \texttt{'a'}, из списка \texttt{['apple',~'banana',~'avocado']}.

    \subsection*{Продвинутый уровень}
    \item Выведите все элементы списка \texttt{[1, 5, 3, 5, 7, 5, 9, 5, 5]}, за исключением числа 5.
    \item Найдите минимальное и максимальное число в списке \texttt{[3, 5, 1, 8, 9, 2]}, не используя функции \texttt{min()} и \texttt{max()}.
    \item Посчитайте количество чисел в списке \texttt{[1, 6, 3, 7, 8, 2, 10]}, которые больше 5.

    % Исправление: добавлена подсказка о структуре решения, так как задача
    % требует отслеживания нескольких переменных состояния одновременно.
    \item Найдите самую длинную последовательность одинаковых символов в строке \texttt{'aabbbccccdddd'}.

    \begin{exercisehint}
        Введите две переменные: \texttt{current} — длина текущей серии одинаковых
        символов, и \texttt{best} — длина наибольшей найденной серии. Обновляйте
        \texttt{best} каждый раз, когда \texttt{current} становится больше.
    \end{exercisehint}

    \item Найдите наибольшее число в списке \texttt{[2, 3, 5, 7, 8, 10, 13]}, которое меньше 10.
    \item Создайте список, содержащий все буквы из строки \texttt{'python programming'}, которые встречаются более одного раза.

    \item Наибольший общий делитель (НОД) нескольких чисел --- это наибольшее целое число, на которое каждое из них делится без остатка.
Например, $\text{НОД}(12, 18, 24) = 6$, так как 6 делит все три числа, а большего такого делителя нет.

Найдите НОД чисел из списка \texttt{[60, 90, 120, 150]} без использования функции \texttt{math.gcd}.

\begin{exercisehint}
    Делитель числа \texttt{n} --- это число \texttt{d}, для которого
    выполняется \texttt{n \% d == 0}.
    Переберите все числа \texttt{d} от 1 до минимального элемента списка:
    большего общего делителя существовать не может, так как НОД обязан
    делить каждый элемент, в том числе наименьший.
    На каждом шаге проверяйте, делится ли на \texttt{d} \emph{каждый}
    элемент списка. Наибольшее подходящее \texttt{d} и будет ответом.
\end{exercisehint}


    \item Найдите все простые числа до 100 методом решета Эратосфена.

    \begin{exercisehint}
        Создайте список из \texttt{True} длиной 101. Начиная с числа 2,
        помечайте все его кратные (4, 6, 8, \ldots) как \texttt{False},
        затем переходите к следующему числу, которое еще не помечено, и повторяйте.
        В конце выведите все индексы, оставшиеся равными \texttt{True}.
    \end{exercisehint}

%     \item Простое число --- это натуральное число больше 1, которое делится без остатка только на 1 и на себя.
% Например: 2, 3, 5, 7, 11, \ldots{} Найдите все простые числа до 100 методом решета Эратосфена.

% \begin{exercisehint}
%     Алгоритм работает со списком булевых значений \texttt{is\_prime}
%     длиной 101, где индекс соответствует проверяемому числу,
%     а значение \texttt{True} означает «число еще не вычеркнуто».
%     Создайте список \texttt{is\_prime = [True] * 101} и пометьте
%     \texttt{is\_prime[0] = is\_prime[1] = False}, так как 0 и 1
%     простыми не являются.
%     Затем для каждого числа \texttt{i} начиная с 2: если
%     \texttt{is\_prime[i]} равно \texttt{True}, пометьте все его
%     кратные \texttt{is\_prime[i*2]}, \texttt{is\_prime[i*3]}, \ldots{}
%     как \texttt{False} --- они составные, так как делятся на \texttt{i}.
%     В конце выведите все индексы \texttt{i},
%     для которых \texttt{is\_prime[i]} осталось равным \texttt{True}.
% \end{exercisehint}

    %\item Найдите все совершенные числа до 1000. (Число называется совершенным, если оно равно сумме всех своих делителей, кроме самого себя.)
    \item Число называется \emph{совершенным}, если оно равно сумме
всех своих делителей, кроме самого себя.
Например: $6 = 1 + 2 + 3$.
Найдите все совершенные числа до 1000.

\begin{exercisehint}
    Делитель числа \texttt{n} --- это число \texttt{d}, для которого
    выполняется \texttt{n \% d == 0}.
    В задаче рассматриваются только \emph{собственные} делители:
    все делители числа \texttt{n}, кроме самого \texttt{n}.
    Заведите переменную-накопитель \texttt{s = 0} и перебирайте
    \texttt{d} от 1 до \texttt{n - 1}: если \texttt{n \% d == 0},
    добавляйте \texttt{d} к \texttt{s}.
    Если в итоге \texttt{s == n} --- число совершенное.
\end{exercisehint}

    %\item Выведите числа Армстронга до 1000. (Число Армстронга — это число, равное сумме своих цифр, возведенных в степень, равную количеству цифр числа. Например: $153 = 1^3 + 5^3 + 3^3$.)
    
    
    \item Число называется \emph{числом Харшад}, если оно делится без остатка
на сумму своих цифр. Например: $18 \div (1 + 8) = 18 \div 9 = 2$ ---
делится, значит 18 является числом Харшад.
Найдите все числа Харшад до 1000.

\begin{exercisehint}
    Алгоритм работает со строковым представлением числа.
    Преобразуйте \texttt{n} в строку \texttt{str(n)}, чтобы получить
    доступ к отдельным цифрам через перебор символов.
    Каждый символ-цифру \texttt{c} преобразуйте обратно в целое число
    \texttt{int(c)} и накапливайте их сумму в переменной \texttt{s}.
    Если \texttt{n \% s == 0} --- число является числом Харшад.
\end{exercisehint}
    % \item Проверьте, является ли число \texttt{25} автоморфным. (Квадрат числа оканчивается самим числом: $25^2 = 625$.)

    % \item Найдите все числа-палиндромы в диапазоне от 1 до 1000, квадрат которых тоже является палиндромом.

    \item Число называется \emph{числом Армстронга}, если оно равно сумме
своих цифр, каждая из которых возведена в степень, равную количеству
цифр числа. Например: $153 = 1^3 + 5^3 + 3^3$.
Выведите все числа Армстронга до 1000.

\begin{exercisehint}
    Алгоритм работает со строковым представлением числа.
    Преобразуйте \texttt{n} в строку \texttt{str(n)},
    чтобы получить доступ к отдельным цифрам через перебор символов
    и узнать их количество \texttt{k = len(str(n))}.
    Каждый символ-цифру \texttt{c} преобразуйте обратно в целое число
    \texttt{int(c)} и возведите в степень \texttt{k}.
    Если сумма всех таких слагаемых равна \texttt{n} ---
    число является числом Армстронга.
\end{exercisehint}

\item Число называется \emph{автоморфным}, если его квадрат
оканчивается на само это число.
Например: $5^2 = 25$;\quad $76^2 = 5776$.
Проверьте, является ли число \texttt{25} автоморфным.

\begin{exercisehint}
    Алгоритм работает со строковыми представлениями числа и его квадрата.
    Количество цифр числа \texttt{n} равно \texttt{k = len(str(n))}.
    Последние \texttt{k} цифр квадрата можно получить срезом
    \texttt{str(n ** 2)[-k:]}: отрицательный индекс отсчитывает
    символы с конца строки.
    Если \texttt{str(n ** 2)[-k:] == str(n)} ---
    число является автоморфным.
\end{exercisehint}

\item Число называется \emph{палиндромом}, если оно читается одинаково
слева направо и справа налево. Например: $12321$, $9009$.
Найдите все числа-палиндромы в диапазоне от 1 до 1000,
квадрат которых тоже является палиндромом.

\begin{exercisehint}
    Преобразуйте целое число в строку: \texttt{s = str(n)}.
    Обращенную копию строки можно получить срезом \texttt{s[::-1]}:
    шаг $-1$ означает перебор символов с конца к началу.
    Если \texttt{s == s[::-1]} --- число является палиндромом.
    Ту же проверку примените к строке \texttt{str(n ** 2)}.
\end{exercisehint}

\end{enumerate}

\end{document}